\documentclass[a4paper,11pt]{article}

\usepackage{amsmath, amssymb, amstext, amsfonts, mathrsfs}


% \sffamily %schrift ohne Serifen

\usepackage[T1]{fontenc} 
% schriftencodierung f�r umlaute, trennung
% f\"ur Uni
\usepackage[latin1]{inputenc}
\usepackage{selinput}
% \usepackage[utf8x]{inputenc} 
\usepackage{bibgerm} 
% german bibliography
\usepackage[german]{babel}
%wichtig f�r deutschen Content
\usepackage{ucs}
%erweiterte UTF-8 Unterst�tzung
\usepackage{wrapfig} 
% Paket zur Positionierung einbinden
\usepackage{multirow}
% zusammenfassen von Tabellenzellen
% \usepackage{subscript}
% zum tiefstellen
\usepackage{lscape}
\usepackage{pdflscape}
% zum drehen der Seite
% \usepackage[super]{natbib}
\usepackage[square,sort,comma,numbers]{natbib}
% Erstellung es Literaturverzeichnisses
\usepackage{url}
% Umbruch f�r URL
\usepackage{pst-3dplot}
% f�r tex Grafiken n�tig
\usepackage{pstricks}
\usepackage{listings}
% f�r das einf�gen von Quelltext
\definecolor{codegray}{rgb}{0.92,0.92,0.92}
\lstset{basicstyle=\fontsize{9}{11}\selectfont\ttfamily, breaklines=true, backgroundcolor=\color{codegray}, numbers=left, numberstyle=\tiny, tabsize=4, language=java}
\definecolor{mymauve}{rgb}{0.58,0,0.82}
\definecolor{mygreen}{rgb}{0,0.6,0}
\lstset{
commentstyle=\color{mygreen},
keywordstyle=\color{mymauve},
language=Java,
stringstyle=\color{blue}
}


%Einstellungen f�r Quellcode
\usepackage[a4paper, left=3cm, right=2cm, top=2cm]{geometry}
% Formatierung R�nder
\usepackage[section]{placeins}
% f�r Floatbarriere
\usepackage{color}
\usepackage{colortbl}
%f�r die Verwendung von Farben

\clubpenalty = 10000 
\widowpenalty = 10000
\displaywidowpenalty = 10000
%Verhinderung von Hurenkindern und Schusterjungen
%10000 bedeutet die sie sollen kommplett vermieden werden

\title{Konzept und prototypische Implementierung eines �bergreifenden Dokumenten- und Medienmanagements als SOA-Anwendung}

\author{Sebastian Rieger}

\pagenumbering{arabic}
%Seitenzahlen(arabische Zahlen)

\setlength{\parindent}{0.25cm} 
%Absatzeinzug �ndern in Zoll
\setlength{\parskip}{0.25cm}
%Absatzabstand
% \linespread {1.1}
%Zeilenabstand

\usepackage{setspace}
\usepackage{hyperref}
%anklickbare Hyperlinks

%funktioniert nicht bei Fu�noten
\usepackage{graphicx}
\usepackage{graphics}
%f�r einbinden von Grafiken

\usepackage{framed}
%f�r Umrandung der Erkl�rung
\usepackage{acronym}
% f�r abk�rzungen
% \usepackage{PSTricks}
\usepackage{epstopdf}
% f�r eps bilder nutze pdflatex --shell-escape this.tex
\usepackage{amssymb}
% f�r mathematische Symbole

\usepackage{hyperref}
% klickbare links

% \usepackage{pdfpages}
\usepackage{rotating}
\usepackage{svg}
%%%%%%%%%%%%%%%%%%%%%%%%%%%%%%%%%%%%%%%%%%%%%%%%%%%%%%%%%%%%%%%%%%%%%%%%%%%%%%%%%%%%%%%%%%%%%%%%%%%%%
%% Angaben zur Arbeit
%%%%%%%%%%%%%%%%%%%%%%%%%%%%%%%%%%%%%%%%%%%%%%%%%%%%%%%%%%%%%%%%%%%%%%%%%%%%%%%

\newcommand{\Autor}{Sebastian Rieger}
\newcommand{\MatrikelNummer}{7406886}
\newcommand{\Kursbezeichnung}{TINF12B1}

\newcommand{\FirmenName}{Karlsruher Institut f�r Technologie (KIT)}
\newcommand{\FirmenStadt}{Karlsruhe}
\newcommand{\FirmenLogoDeckblatt}{{\includegraphics[width=3cm]{Bilder/kitlogo}}}

% Falls es kein Firmenlogo gibt:
%  \newcommand{\FirmenLogoDeckblatt}{}

\newcommand{\BetreuerFirma}{Dipl. -Inform. Thorsten Schlachter}
\newcommand{\BetreuerDHBW}{Dipl. -Inform., MBA Christian Jordan}
\newcommand{\Titel}{Konzept und prototypische Implementierung eines �bergreifenden Dokumenten- und Medienmanagements als SOA-Anwendung}
\newcommand{\AbgabeDatum}{24.08.2015}

\newcommand{\Dauer}{12 Wochen}

% \newcommand{\Abschluss}{Bachelor of Engineering}
\newcommand{\Abschluss}{Bachelor of Science}

\newcommand{\Studiengang}{Angewandte Informatik}
% \newcommand{\Studiengang}{Angewandte Informatik}
\newcommand{\Was}{Bachelorarbeit}

%%%%%%%%%%%%%%%%%%%%%%%%%%%%%%%%%%%%%%%%%%%%%%%%%%%%%%%%%%%%%%%%%%%%%%%%%%%%%%%%%%%%%%%%%%%%%%%%%%%%% 
%steuervariableISO 23081
\usepackage{ifthen} %Package f�r if/else
\newboolean{bilder} %Deklaration
\setboolean{bilder}{true} %Zuweisung
% \setboolean{bilder}{false} %Zuweisung
%%%%%%%%%%%%%%%%%%%%%%%%%%%%%%%%%%%%%%%%%%%%%%%%%%%%%%%%%%%%%%%%%%%%%%%%%%%%%%%%%%%%%%%%%%%%%%%%%%%%%

\begin{document}

\begin{center}
\vspace*{-2cm}
\FirmenLogoDeckblatt\hfill\includegraphics[width=4cm]{Bilder/dhbw-logo}\\[1cm]
{\Huge \Titel}\\[2cm]
{\Huge\scshape \Was}\\[2cm]
{\large f�r die Pr�fung zum}\\[0.5cm]
{\Large \Abschluss}\\[0.5cm]
{\large des Studienganges \Studiengang}\\[0.5cm]
{\large an der}\\[0.5cm]
{\large Dualen Hochschule Baden-W�rttemberg Karlsruhe}\\[0.5cm]
{\large von}\\[0.5cm]
{\large\bfseries \Autor}\\[1cm]
{\large Abgabedatum \AbgabeDatum}
\vfill
\end{center}
\begin{tabular}{l@{\hspace{1cm}}l}
Bearbeitungszeitraum             & \Dauer                       \\
Matrikelnummer                   & \MatrikelNummer              \\
Kurs                             & \Kursbezeichnung             \\
Ausbildungsfirma                 & \FirmenName                  \\
                                 & \FirmenStadt                 \\
Betreuer der Ausbildungsfirma    & \BetreuerFirma               \\
Gutachter der Studienakademie    & \BetreuerDHBW                \\
\end{tabular}

\newpage
%Seitenumbruch
%%%%%%%%%%%%%%%%%%%%%%%%%%%%%%%%%%%%%%%%%%%%%%%%%%%%%%%%%%%%%%%%%%%%%%%%%%%%%%
%% Descr:       Vorlage für Berichte der DHBW-Karlsruhe, Erklärung
%% Author:      Prof. Dr. Jürgen Vollmer, vollmer@dhbw-karlsruhe.de
%% $Id: erklaerung.tex,v 1.2 2010/07/22 13:30:27 vollmer Exp $
%%%%%%%%%%%%%%%%%%%%%%%%%%%%%%%%%%%%%%%%%%%%%%%%%%%%%%%%%%%%%%%%%%%%%%%%%%%%%%%

% In Bachelorarbeiten muss eine schriftliche Erklärung abgegeben werden. In allen anderen
% Arbeiten entf�llt diese. Hierin best�tigen die Studierenden, dass die Bachelorarbeit
% selbst�ndig verfasst und s�mtliche Quellen und Hilfsmittel angegeben sind. Diese Erkl�rung
% bildet das zweite Blatt der Arbeit. Der Text dieser Erkl�rung muss auf einer separaten Seite
% wie unten angegeben lauten.

\newpage
\thispagestyle{empty}
\begin{framed}
\begin{center}
\Large\bfseries Erkl\"arung
\end{center}

\noindent
Gem\"a\ss{}�16 (3) der "`Studien- und Pr\"ufungsordnung f\"ur den Studienbereich
Technik"' vom 1.11.2007.

\medskip
\noindent
Ich habe die vorliegende Arbeit selbstst\"andig verfasst und
keine anderen als die angegebenen Quellen und Hilfsmittel verwendet.

\vspace{3cm}
\noindent
\underline{\hspace{4cm}}\hfill\underline{\hspace{6cm}}\\
Ort~~~~~Datum\hfill Unterschrift\hspace{4cm}
\end{framed}

%%%%%%%%%%%%%%%%%%%%%%%%%%%%%%%%%%%%%%%%%%%%%%%%%%%%%%%%%%%%%%%%%%%%%%%%%%%%%%%
\endinput
%%%%%%%%%%%%%%%%%%%%%%%%%%%%%%%%%%%%%%%%%%%%%%%%%%%%%%%%%%%%%%%%%%%%%%%%%%%%%%%

\newpage
\begin{spacing}{0.9}

%Einf�gen Inhaltsverzeichnis
\tableofcontents



\end{spacing}

\newpage
\section{Abstract}
\newpage
\section{Einleitung} \label{Einfuehrung}
Beh\"orden und Firmen haben heute immer gr\"o\ss{}ere Datenbest\"ande zu verwalten, welche schon in elektronischen Systemen vorzufinden sind. Da diese Systeme historisch bedingt gewachsen sind, entstanden im laufe der Zeit immer mehr kleinere Insell\"osungen, welche ein \"ubergreifendes Verwalten und Nutzen von Daten erschwert oder gar ganz verhindert.

Das Ziel dieser Arbeit soll es daher sein, ein Konzept sowie eine prototypische Implementierung zu ertsellen, welche ein \"ubergreifendes und vereinheitlichtes arbeiten mit den Datenbest\"anden erm\"oglicht. Heutige \ac{DMS} Systeme decken eine vielzahl von Anwendungsgebieten ab, welche an die jeweiligen Problemstellungen wie Rechtsfragen oder fachliche Anforderungen angepasst sind oder angepasst werden k\"onnen.
Ein genaues Konzept mit Anforderungen, soll am Beispiel der \ac{LUBW} sowie der \ac{GAA} entstehen.
\cite{Dokumenten-Management}

Es sollen durch das Konzept auch bestehende Systeme weiter unterst\"utzt werden, wof\"ur hier eine Anbindungsm\"oglichkeit geschaffen werden soll. \"Uber eine serviceorientierte Anbinung sollen somit die verschiedensten Frontends wie Desktopanwendung, Web, Mobil usw. das neue Konzept nutzen k\"onnen. Welches \ac{ECM}-Tool f\"ur die Erstellung des Projekts am besten geeignet ist, muss analysiert werden, ergibt sich jedoch aus den zu verwaltenden Dokumenten.

F\"ur die Arbeit soll hierbei kein komplett neues \ac{DMS} mit Nutzer und Dokumentverwaltung entwickelt werden. Vielmehr soll evaluiert werden, welche Systeme es bereits auf dem Markt gibt, und wie sich diese f\"ur die anstehnde Aufgabe eignen, vor allem im Bezug auf die hierachische Gliederung der Metadaten, wie im Projekt vorgesehen ist. 

Als sogenannter \ac{PoC} soll im Verlauf dieser Arbeit ein bestehendes Fachsystem auf der Basis einer erarbeiteten Schnittstelle nachgebaut werden. Diese Frontend soll das derzeitige \ac{FADO}-System der \ac{LUBW} nachbilden.
\newpage
\section{Lastenheft} \label{Lastenheft}
An das Projekt, welches im Rahmen dieser Arbeit bearbeitet werden soll, gibt es viele Anforderungen. Um alle Anforderungen m\"oglichst strukturiert und \"Ubersichtlich darzustellen, wurde die Form eines Lastenhefts gew\"ahlt.

\subsection{Zielsetzung} \label{Zielsetzung}
Ziel ist es ein einheitliches \ac{ECM} System zu entwerfen, welches m\"oglichst alle Fachdokumente der \ac{LUBW} und der \ac{GAA} beinhaltet.
Hierbei sollen alle bestehenden und zuk\"unfitgen Systeme \"uber eine Schnittstelle mit dem \ac{DMS} kommunizieren.

\subsection{Produkteinsatz} \label{Produkteinsatz}
Bei dem zu entwickelnden System, handelt es sich um einen Prototyp f\"ur ein \ac{DMS}, welches gegebenenfalls von der \ac{LUBW} und der \ac{GAA} eingesetzt werden kann. Das zu entwickelnde Backend, soll alle Dokumente, welche schon heute auf den Websiten zu finden sind zusammenf\"uhren und diese in einer geeigneten Form ablegen.

\subsubsection{Zusammenspiel mit anderen Systemen} \label{Zusammenspiel mit anderen Systemen}
Das Backend soll nicht die heutigen Insell\"osungen f\"ur Dokumentverwaltung der \ac{LUBW} und der \ac{GAA} vereinen. Es sollen sowohl die bestehenden als auch zuk\"unfitge Systeme unterst\"utzt werden. Die jeweils ben\"otigten Daten, sollen hierf\"ur \"uber eine oder gegebenenfalls mehrere Schnittstellen erreichbar sein.

\subsection{Produktfunktionen} \label{Produktfunktionen}
\begin{minipage}{3cm}
/LF10/
\end{minipage}
\begin{minipage}{13cm}
Schon bestehende Dokumente sollen aus den alten Systemen \"ubernommen werden k\"onnen, ohne das ihre Metadaten verloren gehen oder ge\"andert werden. Dies bedeutet, dass das neue System den schon vorhandenen Datenbestand abbilden muss.\\
\end{minipage}
\begin{minipage}{3cm}
/LF20/
\end{minipage}
\begin{minipage}{13cm}
Die Metadaten der einzelnen Dokumente sollen fachlich und technisch gegliedert werden. Hierbei sollen gegebenenfalls Gruppierungen erstellt werden. Zum Beispiel sollen Longitude und Latitude zur Metadaten-Gruppe "`Geodata"' zusammengefasst werden.\\
\end{minipage}
\begin{minipage}{3cm}
/LF30/
\end{minipage}
\begin{minipage}{13cm}
Die Dokumente sollen ebenfalls nach den einzelnen Anwendungen und Fachbereichen gruppiert werden. Hier sollen die Anwendungen auch speziell ihre Dokumente nach Gruppenzugeh\"origkeit abfragen k\"onnen. Dies ist wichtig, da im \ac{FADO} nur die entsprechenden \ac{FADO}-Dokumente angezeigt werden sollen.\\
\end{minipage}
\begin{minipage}{3cm}
/LF40/
\end{minipage}
\begin{minipage}{13cm}
Eine Suche \"uber die Dokumente und ihre Metadaten soll grunds\"atzlich m\"oglich sein. Hier muss der Nutzer aber auch in der Lage sein die Suche weiter einzuschr\"anken oder nur nach bestimmten Metadaten suchen zu k\"onnen.\\
\end{minipage}
\begin{minipage}{3cm}
/LF50/
\end{minipage}
\begin{minipage}{13cm}
Die alten Systeme sollen durch das neue Backend keinerlei Einschr\"ankungen im Funktionsumfang unterworfen sein. Dies muss gegebenenfalls durch verschiedene Arten von Schnittstellen realisert werden.\\
\end{minipage}
\begin{minipage}{3cm}
/LF60/
\end{minipage}
\begin{minipage}{13cm}
F\"ur den Frontend-Prototyp, soll die Kernfunktionalit\"at einer Schnittstelle funktionsf\"ahig sein, anhand dessen der Prototyp evaluiert werden kann.\\
\end{minipage}

\subsection{Produktdaten} \label{Produktdaten}
\begin{minipage}{3cm}
/LD10/
\end{minipage}
\begin{minipage}{13cm}
Die entsprechenden Metadaten zu Dokumenten sollen durch das Backend verwaltet werden. Hierbei sollen nicht nur die manuell hinzugef\"ugten Metadaten beachtet werden, sondern auch die Metadaten, welche das Dateiformat liefert und solche die das Dokument enth\"alt.\\
\end{minipage}
\begin{minipage}{3cm}
/LD20/
\end{minipage}
\begin{minipage}{13cm}
Metadaten sollen nach M\"oglichkeit so zusammengefasst werden, das eine Oberklassenbildung und Vererbung von Metadaten m\"oglich ist. So soll zum Beispiel ein Standort alle Adressdaten zusammenfassen. Ein Standort wiederum kann in mehreren Dokumenten verwendet werden, wie zum Beispiel in Bildern oder in einem Forschungsbericht. Wiederum k\"onnte der Standort aber auch Betandteil der Unterklasse Gericht sein und somit den Standort des Gerichts abbilden.\\
\end{minipage}
\begin{minipage}{3cm}
/LD30/
\end{minipage}
\begin{minipage}{13cm}
Metadaten sollen \"ahnlich wie objektorientierte Klassen in Programmiersprachen agieren. So wird ein Autor "`Max Mustermann"' nur einmal angelegt. Er kann jedoch in mehreren Dokumenten auftauchen. Sucht ein Nutzer nach "`Max Mustermann"', so werden ihm alle Eintr\"age zu diesem angezeigt. Vergleichen kann mal diese Art von Abbildung auch mit einer relationalen Abbildung wie sie in Datenbanken zu finden ist.\\
\end{minipage}
\begin{minipage}{3cm}
/LD40/
\end{minipage}
\begin{minipage}{13cm}
F\"ur die Gliederung der Metdaten sollen auch schon bestehende Standards betrachtet und gegebenenfalls umgesetzt werden.
Im Verlauf der Arbeit muss evaluiert werden, welche Standards es f\"ur Metadaten gibt und wie sie im Projekt umgesetzt werden m\"ussen und k\"onnen.
\cite{Wiki_Dublin_Core} \cite{Wiki_ISO_19115} \cite{Wiki_Exif} \cite{Wiki_Inspire}\\
\end{minipage}


\subsection{Produktleistungen} \label{Produktleistungen}
\begin{minipage}{3cm}
/LL10/
\end{minipage}
\begin{minipage}{13cm}
Das Backend soll f\"ur Dokumente, welche in verschiedenen Versionen vorliegen eine Art Versionsverwaltung bieten, so das auch \"altere Versionen zug\"anglich sind. Dies zum Beispiel bei Gesetzestexten wichtig, da hier immer die zum Vorfall aktuelle Version betrachtet werden muss.\\
\end{minipage}
\begin{minipage}{3cm}
/LL20/
\end{minipage}
\begin{minipage}{13cm}
Die Dokumente sollen \"uber verschiedene Schnittstellen wie zum Beispiel eine \ac{REST}-Schnittstelle abgerufen werden k\"onnen. Der Einsatz von \ac{API}s ist je nach \ac{ECM}-Tool gegebenenfalls auch m\"oglich, jedoch ist eine solche Ums\"atzung nicht Teil der Arbeit. \cite{Wiki_REST}\\
\end{minipage}
\begin{minipage}{3cm}
/LL30/
\end{minipage}
\begin{minipage}{13cm}
Eine Suche \"uber Metadaten muss realisert werden. Hierbei soll dem Nutzer zum einen eine grobe aber auch zum anderen eine verfeinerte eingeschr\"ankte Suche nach bestimmten Metadaten geboten werden.\\
\end{minipage}
\section{Analyse Datenbest\"ande}
\subsection{FADO und Untergruppe}
\subsection{DRS}
\subsection{Bilddatenbank}
\section{Stand der Technik / existierende Konzepte}
\subsection{Fachlich}
\subsection{Technisch}
\section{Erstellung eines Datenkonzepts}
\subsection{Fachlich}
\subsection{Technisch}
\section{Technologievergleich}
\subsection{Alfresco}
\subsection{anderes ECM}
\subsection{anderes ECM}
\section{Implementierung}
\subsection{Schnittstelle}
\subsection{Backend}
\subsection{Frontend}


\newpage
\section{Abk�rzungsverzeichnis}
\begin{acronym}
%   \acro{SDK}{\emph{Software Development Kits}}
%   \acro{OHA}{\emph{Open Handset Alliance}}
%   \acro{IDE}{\emph{Integrierte Entwicklungsumgebung}}
%   \acro{ADT}{\emph{Android Development Tools}}
%   \acro{DVM}{\emph{Dalvik Virtual Machine}}
%   \acro{FME}{\emph{Funkmeldeempf\"anger}}
%   \acro{JIT}{\emph{Just in Time}}
%   \acro{ART}{\emph{Android Runtime}}
%   \acro{AOT}{\emph{Ahead-of-Time-Decodierung}}
%   \acro{API}{\emph{Application Programming Interface}}
  \acro{ECM}{\emph{Enterprise Content Management}}
  \acro{DMS}{\emph{Dokumentenmanagementsystem}}
  \acro{PoC}{\emph{Proof of Concept}}
  \acro{LUBW}{\emph{Landesanstalt f\"ur Umwelt, Messungen und Naturschutz Baden-W\"urttemberg}}
  \acro{GAA}{\emph{Gewerbeaufsicht Baden-W\"urttemberg}}
  \acro{REST}{\emph{Representational State Transfer}}
  \acro{FADO}{\emph{Fachdokumente Online}}
  \acro{DRS}{\emph{Document Retrieval System}}
  \acro{INSPIRE}{\emph{Infrastructure for Spatial Information in the European Community}}
  \acro{API}{\emph{Application Programming Interface}}
  \acro{EU}{\emph{Europ\"aische Union}}
  \acro{}{\emph{}}
%  \acro{KIT}{\emph{Karlsruher Instituts f�r Technologie}}
%  \acro{GDS}{\emph{Generic Data Services}}
% %  \acro{LSDF}{\emph{Large Scale Data Facility}}
%  \acro{OPM}{\emph{Objektorientierten Programmiermodell}}
%  \acro{SMD}{\emph{Strukturelle Metadaten}}
%  \acro{JSON}{\emph{JavaScript Object Notation}}
% %  \acro{HALO}{\emph{High Altitude and Long Range Research Aircraft}}
%  \acro{IAI}{\emph{Institut f�r Angewandte Informatik}}
%  \acro{JAXB}{\emph{Java Architecture for XML Binding}}
%  \acro{UDDE}{\emph{User Data Description Editor}}
%  \acro{AMD}{\emph{Anwendermetadaten}}
%  \acro{CG}{\emph{Class Generator}}
%  \acro{IG}{\emph{Interface Generator}}
\end{acronym}
\newpage
\listoffigures
\newpage
\listoftables
% Abk�rzungsverzeichnis
\newpage
\bibliographystyle{alpha}
% verzeichnis im DIN format
\bibliography{Quellen}
\end{document}
